\chapter{Pendahuluan}
Pada bab ini kita akan sedikit membahas tentang sejarah dari AI.
Bagi yang tidak sabar dapat langsung loncat kepada bagian yang
mungkin lebih menarik, tetapi saya sarankan suatu saat baca juga
bagian ini agar dapat lebih memahami esensinya.

\section{Apa itu Kecerdasan}
Apa itu \textbf{kecerdasan} (\emph{intelligence})?
Ini adalah pertanyaan yang dapat dilihat dari berbagai kacamata;
filsafat, psikologi, neuroscience, dan seterusnya.
Ada beberapa definisinya, antara lain:
\begin{itemize}
\item kemampuan umum untuk memahami, belajar, bernalar, memecahkan masalah, 
   beradaptasi dengan situasi baru, serta menggunakan pengetahuan dan 
   pengalaman secara efektif dalam menghadapi tantangan kehidupan;
\item sekumpulan kemampuan kognitif yang memungkinkan makhluk hidup 
   atau sistem buatan untuk belajar, memahami, beradaptasi, dan 
   bertindak secara purposif dalam lingkungannya;
\item kemampuan untuk memahami, belajar, berpikir, memecahkan masalah, 
   dan beradaptasi dengan lingkungan baru.
\end{itemize}

Pada jaman dahulu orang yang cerdas adalah yang mahir di bidang keilmuan
tertentu, seperti misalnya matematika. Yang pandai berhitung adalah
orang yang cerdas.
Atau bisa juga orang yang sangat pandai bermain catur (atau Go) adalah
orang yang cerdas.
Maka untuk membuat sebuah alat atau mesin yang cerdas adalah membuat
sebuah mesin yang bisa bermain catur (dan mengalahkan manusia).
Diceritakan ada sebuah mesin yang katanya bisa bermain catur dengan baik.
Salah satu mesin ini dikenal dengan nama \emph{``the Mechanical Turk''}.
Ternyata mesin hebat ini adalah bohongan karena sebetulnya di dalam
sebuah kotak besar (mesin) ini ada orang yang berukuran mini.
Jadi sebetulnya isinya orang juga. Bukan mesin.
Obsesi manusia dalam membuat mesin yang pandai memainkan catur ini
akhirnya berhasil. IBM dengan mesin \emph{Deep Blue}-nya akhirnya
mengalahkan juara catur dunia, Garry Kasparov, di tahun 1997.

Pada kesempatan yang lain, Google dengan mesinnya \emph{AlphaGo}
yang dikembangkan oleh tim Google DeepMind
akhirnya mengalahkan manusia dalam permaian Go.
Permainan Go sendiri seringkali dianggap lebih kompleks dibandingkan
dengan permainan catur.

Meskipun manusia sudah dapat membuat mesin yang ``cerdas'', dalam artian
dapat mengalahkan manusia dalam catur dan Go, tetapi sifat kecerdasan
ini masih terbatas. Mesin tersebut belum dapat kita gunakan untuk
berdiskusi, misalnya. Kehebatan manusia dalam membuat mesin cerdas
masih terbatas. Mesin tersebut belum bisa ``sadar'' \emph{(conscious)}
seperti lazimnya seorang manusia.
Dapatkah kita membuat mesin yang sadar?
Nanti ini kita bahas di bagian-bagian selanjutnya.

Kecerdasan ini seringkali diukur dengan menggunakan IQ (Intelligence
Quotient) untuk alasan praktis. Ada sebuah angka yang terukur.
Namun masih banyak perdebatan tentang ukuran ini.


\section{Perioda Awal: 1950-1980}
Istilah \emph{``artificial intelligence''} dicetuskan (di-koin-kan) oleh
John McCarthy di tahun 1956 ketika dia menuliskan proposal untuk
\emph{1956 Dartmouth Summer Research Project on Artificial Intelligence}.
(Proposal tersebut juga ditandatangani oleh
Marvin Minsky, Nathaniel Rochester, dan Claude Shannon.
Mereka termasuk orang-orang yang terkenal di bidang ilmu komputer.)

\begin{quote}
   \emph{The study is to proceed on the basis of the conjecture that every aspect 
   of learning or any other feature of intelligence can in principle be so 
   precisely described that a machine can be made to simulate it.}
\end{quote}

Sejak saat ini ada banyak penelitian terkait dengan AI.
Pemerintah Amerika Serikat mulai mengucurkan dana untuk penelitian AI.

Pada tahun 1966 ada sebuah program yang melakukan tanya jawab dengan
pengunanya secara interaktif yang bernama Eliza. Seperti \emph{ChatGPT}.
Kalau sekarang, program ini mungkin disebut sebagai \emph{chatbot}.
Sebetulnya program ini tidak memiliki kecerdasan sebagaimana dimasudkan
dengan AI, tetapi berkesan cerdas.
Jawaban yang diberikan oleh program Eliza tersebut sebetulnya sudah ada
dalam bentuk daftar (list) atau pilihan. Program tinggal memilih Jawaban
yang paling mendekati dengan pertanyaan tersebut.
Saya menuliskan ulang program ini dalam bahasa BASIC ketika pertama kali
saya belajar pemrograman di tahun 1980.

Perlu diingat bahwa pada masa itu yang disebut komputer baru berupa
komputer mainframe yang jumlahnya sedikit dan kemampuan komputasinya
sangat terbatas.
Demikian pula rangkaian elektronik masih terbatas, belum digital,
sehingga solusi-solusi elektronika juga masih terbatas.
Akibat keterbatasan ini kompleksitas masalah yang dapat dipecahkan oleh AI
sangat terbatas.
Ada satu masalah yang disebut \emph{``XOR problem''} yang tidak dapat
dipecahkan oleh AI (pada saat itu).
Maka AI dikatakan penipuan (fraud) sehingga dana penelitian pun banyak
yang dihentikan. \emph{AI winter.}


\section{Perioda: 1980-2010}